\documentclass[12pt]{article}
\usepackage[speech=prose, teiexport=tidy]{ekddrama}

\SetLineation{
    modulo
}

\title{MONSIEUR DE POURCEAUGNAC}
\author{Molière}
\date{1665}
%% https://theatre-classique.fr/pages/pdf/MOLIERE_MONSIEURPOURCEAUGNAC.pdf %%

\DeclareWitness{A}{A}{\emph{BnF} RES-YF-3207}[
               settlement=Paris,
               idno=RES-YF-3207,
               institution=Bibliothéque Nationale de France,
               origDate=1669]
\DeclareWitness{B}{B}{\emph{BnF} RES8ERVE RF-4597}[
               settlement=Paris,
               idno=RESERVE 8-RF-4597,
               institution=Bibliothéque Nationale de France]
\DeclareWitness{C}{C}{\emph{BnF} RESERVE 8-RF-4596}[
               settlement=Paris,
               idno=RESERVE 8-RF-4596,
               institution=Bibliothéque Nationale de France,
               origDate=1669]
\DeclareShorthand{a}{α}{A,B,C}


% Define speakers
\DeclareCast{julie}{JULIE}{fille d'Oronte}
\DeclareCast{eraste}{ERASTE}{amant de Julie}
\DeclareCast{nerine}{NÉRINE}{femme d'intrigue}

\begin{document}


\begin{xltabular}[c]{0.75\linewidth}{lXl}
  \caption*{\textbf{Conspectus Siglorum}\label{tab:conspectus-siglorum}}\\
  \SigLine{A}\\
  \SigLine{B}\\
  \SigLine{C}\\
\end{xltabular}

\newpage

\begin{xltabular}{0.75\linewidth}{l X l}
  \caption*{\textbf{Cast List}\label{tab:cast-list}}\\
  \hline
  \textbf{Personnages} & \textbf{Description} & \textbf{Acteur} \\
  \hline
  \CastTableBody
  \hline
\end{xltabular}


\newpage

\begin{ekddrama}
\ekdscene{head=SCÈNE I., n=1, type=scene}
\begin{speech}
%Julie, Eraste, Nérine.
\loc{julie}
Mon Dieu ! Eraste, gardons d'être surpris ; je tremble
qu'on ne nous voie ensemble, et tout serait \app{\lem[wit=A]{perdu}\rdg[wit=C]{perdue}}, après la
défense que l'on m'a faite.
\loc{eraste}
Je regarde de tous côtés, et je n'aperçois rien.
\loc{julie}
Aie aussi l'oeil au guet, Nérine, et prends bien garde qu'il
ne vienne personne.
\loc{nerine}
Reposez-vous sur moi, et dites hardiment ce que vous
avez à vous dire.
\loc{julie}
Avez-vous imaginé pour notre affaire quelque chose de
favorable ? Et croyez-vous, Eraste, pouvoir venir à bout
de détourner ce fâcheux mariage que mon père s'est mis
en tête ?
\loc{eraste}
Au moins y travaillons-nous fortement ; et déjà nous
avons préparé un bon nombre de batteries pour renverser
ce dessein ridicule.
\loc{nerine}
Par ma foi ! Voilà votre père.
\loc{julie}
Ah ! Séparons-nous vite.
\loc{nerine}
Non, non, non, ne bougez : je m'étais trompée.
\loc{julie}
Mon Dieu ! Nérine, que tu es sotte de nous donner de ces
frayeurs !
\loc{eraste}
Oui, belle Julie, nous avons dressé pour cela quantités de
machines, et nous ne feignons point de mettre tout en
usage, sur la permission que vous m'avez donnée. Ne
nous demandez point tous les ressorts que nous ferons
jouer : vous en aurez le divertissement ; et, comme aux
comédies, il est bon de vous laisser le plaisir de la
surprise, et de ne vous avertir point de tout ce qu'on vous
fera voir. C'est assez de vous dire que nous avons en
main divers stratagèmes tous prêts à produire dans
l'occasion, et que l'ingénieuse Nérine et l'adroit Sbrigani
entreprennent l'affaire.
\loc{nerine}
Assurément. Votre père se moque-t-il de vouloir vous
anger de son avocat de Limoges, Monsieur de
Pourceaugnac, qu'il n'a vu de sa vie, et qui vient par le
coche vous enlever à notre barbe ? Faut-il que trois ou
quatre mille écus de plus, sur la parole de votre oncle, lui
fassent rejeter un amant qui vous agrée ? Et une personne
comme vous est-elle faite pour un Limousin ? 

S'il a envie de se marier, que ne prend-il une Limosine et ne
laisse-t-il en repos les chrétiens ? Le seul nom de
Monsieur de Pourceaugnac m'a mis dans une colère
effroyable. J'enrage de Monsieur de Pourceaugnac.
Quand il n'y aurait que ce nom-là, Monsieur de
Pourceaugnac, j'y brûlerai mes livres, ou je romprai ce
mariage, et vous ne serez point Madame de
Pourceaugnac. Pourceaugnac ! Cela se peut-il souffrir ?
Non, Pourceaugnac est une chose que je ne saurais
supporter ; et nous lui jouerons tant de pièces, nous lui
ferons tant de niches sur niches ; que nous renverrons à
Limoges Monsieur de Pourceaugnac.
\loc{eraste}
Voici notre subtil Napolitain, qui nous dira des nouvelles.
\end{speech}

\end{ekddrama}
\end{document}